\chapter{Chapter Two}

Start writing chapter two here.% % Detailed LaTeX document for Lab 2: Enhancing the Agent with Memory
% % This standalone chapter explains how to add short‑term and long‑term memory
% % to your customer‑support agent using Amazon Bedrock AgentCore Memory and
% % Strands Agents.  It is self‑contained and includes code listings and
% % discussions.  The narrative is informed by the AWS workshop materials,
% % but all special characters have been removed for compatibility with LaTeX.

% \documentclass[12pt]{article}

% % Packages for layout and formatting
% \usepackage{geometry}
% \geometry{margin=1in}
% \usepackage{setspace}
% \setstretch{1.2}
% \usepackage{amsmath, amssymb}
% \usepackage{graphicx}
% \usepackage{hyperref}
% \hypersetup{hidelinks}
% \usepackage{xcolor}
% \usepackage{listings}

% % Define colours for code listings
\definecolor{keyword}{RGB}{0,0,128}
\definecolor{comment}{RGB}{0,96,0}
\definecolor{string}{RGB}{128,0,0}
\lstset{
  language=Python,
  basicstyle=\ttfamily\small,
  keywordstyle=\color{keyword}\bfseries,
  commentstyle=\color{comment}\itshape,
  stringstyle=\color{string},
  showstringspaces=false,
  frame=single,
  framerule=0.5pt,
  numbers=left,
  numberstyle=\tiny,
  numbersep=5pt
}

% \begin{document}

% \title{Lab 2: Enhancing the Customer‑Support Agent with Memory}
% \author{}
% \date{}
% \maketitle

% \tableofcontents

\chapter{Lab 2: Enhancing the Customer‑Support Agent with Memory}
\label{chap:Lab 2: Enhancing the Customer‑Support Agent with Memory}
\section{Introduction and Goals}

\section{Introduction and Motivation}

In Lab 1 you created a prototype customer‑support agent that could answer
questions using specialised tools and a knowledge base.  However, the agent
was entirely stateless: every conversation started from scratch.  If a
customer returned with a follow‑up question a week later, the agent had no
recollection of previous interactions.  This leads to repetitive
conversations, frustrated users and missed opportunities for personalisation.

\emph{Amazon Bedrock AgentCore Memory} solves this problem by providing a
managed service that allows agents to remember context across sessions.
AgentCore Memory supports both short‑term memory (for the duration of a
conversation) and long‑term memory (persisted across sessions).  By the end
of this lab your agent will:

\begin{itemize}
  \item Capture and persist key pieces of conversation state, such as user
        preferences and factual information.
  \item Recall preferences and facts when the same customer returns,
        providing personalised responses.
  \item Seamlessly update the memory after each interaction so the agent
        continually learns.
\end{itemize}

This chapter explains the underlying concepts, shows how to create and
configure memory resources, demonstrates how to seed memory with sample
interactions and integrates memory into the agent application.

\section{AgentCore Memory Concepts}

AgentCore Memory distinguishes between short‑term memory (STM) and
long‑term memory (LTM):

\begin{itemize}
  \item \textbf{Short‑Term Memory (STM)} persists within a single
        conversation or closely related sessions.  It captures the dialogue
        history and helps the agent maintain coherence from one message to
        the next.
  \item \textbf{Long‑Term Memory (LTM)} stores facts and preferences
        extracted from conversations over time.  It allows the agent to
        remember customers across sessions, recall preferences and build
        persistent profiles.
\end{itemize}

Long‑term memory is further organised into \emph{strategies}.  Each
strategy defines how information is extracted and stored.  The workshop
uses two built‑in strategies:

\begin{description}
  \item[User Preference] Learns customers’ preferences and behavioural
    patterns (for example, preferred notification channels or favourite
    product categories).
  \item[Semantic] Stores factual information about previous conversations and
    orders (for example, ``customer reported screen flickering on laptop'').
\end{description}

Memories are saved under \emph{namespace} patterns.  A namespace is a path
that groups related memory entries.  The lab uses two namespaces:
\begin{center}
\begin{tabular}{ll}
\textbf{Namespace} & \textbf{Purpose} \\\hline
\texttt{support/customer/\{actorId\}/preferences} & stores user preference
records for each customer \\
\texttt{support/customer/\{actorId\}/semantic} & stores semantic facts and
context for each customer
\end{tabular}
\end{center}

The curly braces \texttt{\{actorId\}} are placeholders that will be
substituted with the actual user identifier at runtime.  This grouping
ensures that each customer’s memory is isolated from others.

\section{Creating and Configuring a Memory Resource}

Before your agent can use AgentCore Memory, you must create a memory
resource.  In a production system this is typically done once via the AWS
console or a provisioning script.  For the purposes of this lab, you can
create the resource programmatically using the \texttt{MemoryClient}
class.  The following example creates a memory with both user‑preference
and semantic strategies and exports the generated memory ID:

\begin{lstlisting}[language=Python]
import os
from bedrock_agentcore.memory import MemoryClient

# Initialize a client in your region
client = MemoryClient(region_name="us-east-1")

# Create a comprehensive memory resource with both strategies
memory_resource = client.create_memory_and_wait(
    name="SupportAgentMemory",
    description="Memory with user preferences and semantic facts",
    strategies=[
        {
            "userPreferenceMemoryStrategy": {
                "name": "PreferenceLearner",
                "namespaces": ["support/customer/{actorId}/preferences"]
            }
        },
        {
            "semanticMemoryStrategy": {
                "name": "FactExtractor",
                "namespaces": ["support/customer/{actorId}/semantic"]
            }
        }
    ]
)

# Save the memory ID for later use (e.g. to an environment variable)
memory_id = memory_resource.get('id')
print(f"Created memory with ID: {memory_id}")
os.environ['AGENTCORE_MEMORY_ID'] = memory_id
\end{lstlisting}

The call to \texttt{create\_memory\_and\_wait} blocks until the memory
resource is fully provisioned, which may take several seconds.  If you
prefer not to wait, you can call \texttt{create\_memory} and poll the
status manually.  Once created, store the memory ID securely (for
example, in AWS Systems Manager Parameter Store) so your agent can use
it later.

\section{Seeding Memory with Sample Interactions}

To demonstrate long‑term memory, it helps to pre‑populate the resource
with a few sample events.  Each event represents a previous customer
interaction.  In production these events would be generated automatically
by your agent, but in this lab you will create them manually.  An event
includes the following fields:

\begin{itemize}
  \item \texttt{memory\_id}: the ID of the memory resource.
  \item \texttt{actor\_id}: a unique identifier for the user (e.g.
        \texttt{customer\_001}).
  \item \texttt{session\_id}: a unique identifier for the session during
        which the event occurred.
  \item \texttt{timestamp}: when the event happened.
  \item \texttt{type}: either \texttt{USER\_PREFERENCE} or
        \texttt{SEMANTIC}, indicating which strategy should handle the
        event.
  \item \texttt{payload}: a dictionary containing the actual content of the
        memory.
\end{itemize}

Here is a helper function to create preference and semantic events for a
single customer.  It uses the \texttt{create\_event} method of
\texttt{MemoryClient}:

\begin{lstlisting}[language=Python]
from datetime import datetime
from uuid import uuid4

# Assume `client` is the MemoryClient from the previous section

def seed_customer_memory(memory_id: str, actor_id: str):
    """Seed user preference and semantic memory events for a customer."""
    session_id = f"seed_session_{uuid4().hex}"
    timestamp = datetime.utcnow().isoformat()

    # Define a user preference event: the customer prefers email updates
    preference_event = {
        "memory_id": memory_id,
        "actor_id": actor_id,
        "session_id": session_id,
        "timestamp": timestamp,
        "type": "USER_PREFERENCE",
        "payload": {
            "channel": "email",
            "newsletter_opt_in": True,
            "preferred_device": "laptop"
        },
    }

    # Define a semantic event: factual information about a previous issue
    semantic_event = {
        "memory_id": memory_id,
        "actor_id": actor_id,
        "session_id": session_id,
        "timestamp": timestamp,
        "type": "SEMANTIC",
        "payload": {
            "issue": "battery draining quickly",
            "product": "smartphone",
            "troubleshooting_steps": [
                "reduce screen brightness",
                "disable background apps",
                "update firmware"
            ]
        },
    }

    # Create the events via the API
    client.create_event(**preference_event)
    client.create_event(**semantic_event)
    print(f"Seeded memory for {actor_id}")

# Example usage
seed_customer_memory(memory_id, actor_id="customer_001")
\end{lstlisting}

After running this helper, the memory resource will contain two events
associated with the user \texttt{customer\_001}: one storing their
preferences and another storing a factual description of a previous
problem.  You can repeat this process for additional customers or
sessions.

\section{Retrieving Memory and Configuring Retrieval}

When your agent starts a new session, it needs to know which pieces of
memory to retrieve.  The \texttt{AgentCoreMemoryConfig} class defines
both the identifiers (memory ID, actor ID and session ID) and the
retrieval policy.  You can specify retrieval behaviour globally for all
namespaces or fine‑tune it per namespace using \texttt{RetrievalConfig}.

For example, the following configuration retrieves up to three user
preferences and five semantic facts for the current actor:

\begin{lstlisting}[language=Python]
from datetime import datetime
from bedrock_agentcore.memory.integrations.strands.config import (
    AgentCoreMemoryConfig,
    RetrievalConfig,
)

# Assume the memory ID is stored in an environment variable
MEMORY_ID = os.environ['AGENTCORE_MEMORY_ID']

# Generate identifiers for this session
ACTOR_ID = "customer_001"
SESSION_ID = f"session_{uuid4().hex}"

# Configure retrieval: different top_k values per namespace
retrieval_policy = {
    "support/customer/{actorId}/preferences": RetrievalConfig(
        top_k=3,
        relevance_score=0.5,
    ),
    "support/customer/{actorId}/semantic": RetrievalConfig(
        top_k=5,
        relevance_score=0.3,
    ),
}

memory_config = AgentCoreMemoryConfig(
    memory_id=MEMORY_ID,
    actor_id=ACTOR_ID,
    session_id=SESSION_ID,
    retrieval_config=retrieval_policy,
)
\end{lstlisting}

The \texttt{relevance\_score} parameter (between 0 and 1) controls how
strictly the retrieval engine ranks memories by relevance.  Lower values
will include more distant memories; higher values require very similar
content.

\section{Integrating Memory into the Agent}

To enable memory within the agent, you wrap the memory configuration in a
session manager.  The Strands Agents SDK provides
\texttt{AgentCoreMemorySessionManager}, which automatically retrieves
memories before each invocation and persists new events afterwards.

Here is how to create a session manager and pass it to the agent along
with your tools and model.  Notice that we keep the same system prompt
and tools from Lab 1 but now supply the session manager:

\begin{lstlisting}[language=Python]
from bedrock_agentcore.memory.integrations.strands.session_manager import (
    AgentCoreMemorySessionManager,
)
from strands import Agent

# Create the session manager for this user and session
session_manager = AgentCoreMemorySessionManager(
    agentcore_memory_config=memory_config,
    region_name="us-east-1",
)

# Create the agent with memory
agent_with_memory = Agent(
    system_prompt=SYSTEM_PROMPT,
    tools=[get_return_policy, get_product_info, web_search, get_technical_support],
    model=model,
    session_manager=session_manager,
)

# Invoke the agent: it will automatically retrieve relevant memories
# before answering and store new events afterwards
response = agent_with_memory.invoke(
    "I recently bought a smartphone; what do you recommend to improve battery life?"
)
print(response)
\end{lstlisting}

When you call \texttt{invoke}, the session manager performs the following
workflow automatically:

\begin{enumerate}
  \item It retrieves relevant memories from the configured namespaces and
        passes them to the language model as part of the context.  This
        ensures the model is aware of the user’s preferences and past
        issues.
  \item After the agent generates a response, the session manager
        extracts new long‑term memory messages (for example, updated
        preferences) and saves them back to AgentCore Memory via the
        \texttt{create\_event} API.  Short‑term memory for the current
        conversation is maintained automatically.
\end{enumerate}

By externalising memory handling into a session manager, you keep your
agent code clean and focused on business logic.  The memory manager
handles retrieval and persistence behind the scenes.

\section{Testing the Memory‑Enabled Agent}

Let’s simulate a returning customer.  In the first interaction the
customer tells the agent they prefer email updates and reports a
battery‑drain issue on their smartphone.  The agent seeds the memory
accordingly.  In a subsequent session, the agent recalls this information
and personalises its response.

\begin{lstlisting}[language=Python]
# First interaction: set user preference and report issue
first_session_id = f"session_{uuid4().hex}"
memory_config.actor_id = "customer_002"
memory_config.session_id = first_session_id

# Create a temporary session manager
first_session_manager = AgentCoreMemorySessionManager(memory_config, region_name="us-east-1")
agent_temp = Agent(
    system_prompt=SYSTEM_PROMPT,
    tools=[get_return_policy, get_product_info, web_search, get_technical_support],
    model=model,
    session_manager=first_session_manager,
)

# Customer expresses their preference in the conversation
agent_temp.invoke("I prefer to receive updates by email.")
# Customer reports a technical problem
agent_temp.invoke("My smartphone battery is draining quickly.")

# Second interaction: same customer returns later
second_session_id = f"session_{uuid4().hex}"
memory_config.session_id = second_session_id
second_session_manager = AgentCoreMemorySessionManager(memory_config, region_name="us-east-1")
agent_memory = Agent(
    system_prompt=SYSTEM_PROMPT,
    tools=[get_return_policy, get_product_info, web_search, get_technical_support],
    model=model,
    session_manager=second_session_manager,
)

# Ask a follow‑up question; the agent should recall the preference and issue
response = agent_memory.invoke("Do you remember how I like to get updates and my last issue?")
print(response)
\end{lstlisting}

You should see that the agent greets the customer as a returning user,
recalls their preferred contact method (email) and summarises the
previous battery‑drain issue.  This demonstrates that AgentCore Memory
has successfully persisted and retrieved long‑term knowledge.

\section{Hooking into Memory with Custom Hooks (Optional)}

Strands Agents provides a powerful hook system that allows you to
intercept various events in the agent lifecycle.  You can use hooks to
modify the prompt, inject custom context, or perform side effects like
logging.  In the memory lab, hooks are used to automatically transform
agent outputs into long‑term memory events without writing explicit
\texttt{create\_event} calls.

To implement hooks you extend the \texttt{HookProvider} class and
override callback methods such as \texttt{before\_invoke} and
\texttt{after\_invoke}.  The workshop’s \texttt{CustomerSupportMemoryHooks}
class extracts preferences from messages and stores them.  While the full
implementation is beyond the scope of this document, consult the workshop
notebook for a complete example.  Hooks are an advanced feature that can
further automate memory management and customise retrieval behaviour.

\section{Limitations and Next Steps}

Adding memory significantly improves the user experience: the agent
recognises returning customers, remembers preferences and provides
context‑aware answers.  Nevertheless, there remain areas for
improvement:

\begin{itemize}
  \item \textbf{Identity and access control.}  Memory is associated with an
        \emph{actor ID} provided by your application.  In Lab 3 you will
        integrate authentication so that only authorised users can access
        their own data and enterprise APIs.
  \item \textbf{Shared tool access.}  Currently all tools are defined in
        your code.  Lab 3 introduces the AgentCore Gateway to expose tools
        as microservices and to share them across multiple agents.
  \item \textbf{Deploying to production.}  In Lab 4 you will wrap your
        memory‑enabled agent in a serverless runtime on Bedrock
        AgentCore, enabling auto‑scaling and observability.
  \item \textbf{Evaluations and UI integration.}  Labs 5 and 6 add online
        evaluation metrics and a user interface, respectively, to monitor
        quality and deliver a polished experience.
\end{itemize}

By completing this lab you have transformed your simple prototype into a
context‑aware assistant that builds relationships over time.  The
subsequent labs will build on this foundation to deliver a secure,
scalable and user‑friendly solution.



